% Einheit 1 von 2 – Der Begriff (bank/uneigentliche-integrale/e1.jsonl, zusammenbau v0.1)
%% TODO Zeitmarke Einheit 1: Marken-Zeile nicht lesbar
%% TODO Zweigzeile Teil 1: was hier gelernt wird (Satz in Schülersprache aus dem Einheitstitel) – Einheit 1
\einheitenkopf[e1]{Einheit 1 von 2 · Der Begriff}
\zweigzeile{keine Prüfungsaufgabe}
\setcounter{aufgabe}{5}

%% TODO Titel: Ich-kann-Satz fehlt in der Bank (Platzhalter: Kettenname) – Nr. 6
%% TODO Anweisung: ein Satz über der ersten Teilaufgabe fehlt in der Bank – Nr. 6
\begin{aufgabe}{Der Begriff}
%% TODO \rechenplatz{n}: Schrittzahl des Grundfalls fehlt in der Bank (nur bei fester Schreibform, 2.3 b)
\begin{teile}
\teil Der Graph von $f$ ist gezeichnet. Die Fläche zwischen Graph und x-Achse rechts der y-Achse – endet sie? \\ \kreuz{endet an einer Grenze} \\ \kreuz{läuft ohne Ende aus}

\begin{ksys}[xmin=0,xmax=8,ymin=0,ymax=4]\funktion{-0.5*\x+2}{f}\end{ksys}
\teil $f(x) = e^{-2x}$ – Inhalt $A(w)$ der Fläche zwischen dem Graphen von $f$ und der x-Achse von $0$ bis $w$ als Term in $w$? A(w) = \leerfeld
\teil $f(x) = 2e^{-0{,}5x}$ – Inhalt der unbegrenzten Fläche zwischen Graph und x-Achse rechts der y-Achse?
\teil Vergleiche $∫_1^w \frac{1}{x^2}\,dx$ und $∫_1^w \frac{1}{x}\,dx$ für $w \to \infty$: Welche der beiden unbegrenzten Flächen hat einen endlichen Inhalt?
\end{teile}
\end{aufgabe}

%% TODO Titel: Ich-kann-Satz fehlt in der Bank (Platzhalter: Kettenname) – Nr. 7
%% TODO Anweisung: ein Satz über der ersten Teilaufgabe fehlt in der Bank – Nr. 7
\begin{aufgabe}{Der Begriff – Fehler finden, Begründen}
\begin{teile}
\teil Ich finde den Fehler: Zu zeigen ist, dass $A(p) = 6 - 6e^{-p}$ stets kleiner als $6$ ist. Lösung: „Für $p \to \infty$ geht $A(p)$ gegen $6$, also ist $A(p) < 6$.“
\teil Begründe: $A(p) = 10 - 10e^{-0{,}1p}$ bleibt für jedes $p$ unter $10$.
\end{teile}
\end{aufgabe}
